\chapter{おわりに}
\thispagestyle{myheadings}

\section{まとめ}
本研究は，近年ハイブリッドワークや柔軟な勤務制度の普及により，出社の有無や在室時間が個人ごとに分散し，従来のような偶発的な対面コミュニケーションが生じにくくなっている状況を背景とし，インフォーマルなコミュニケーションの支援を目的とした．
出社頻度や滞在時間が個人ごとに分散する環境では，対面でのコミュニケーションの機会が十分に活用されにくい．
特に退室時はインフォーマルなコミュニケーションが生じやすい状況であるにもかかわらず，このタイミングを対象とした支援は必ずしも十分に検討されていない．
そこで本研究では，帰宅予測時刻に基づいた複数の滞在者が共に帰宅できそうな時刻を提示するシステム「leaving-match」を提案した．
本システムでは，ユーザ自身が声をかけなくても同行可能なタイミングを把握できるようにし，心理的負担の軽減を通じて退室時の同行とそれに伴うインフォーマルなコミュニケーションの発生を目指す．


提案システムでは複数の滞在者の帰宅予測時刻から一つの時刻を算出し，バスの時刻と組み合わせた提示をする．
帰宅推奨時刻の算出には，まず先行研究から滞在者の帰宅予測時刻を取得し，人数が最大となる区間の抽出，距離が最短の点を中心に距離で重み付けする．
算出された帰宅推奨時刻に対して，その時刻に最も近いバスの発車時刻を基準とし，さらにその前後に位置するバスの発車時刻を含めた三つの時刻を取得し提示を行い，投票を通じて一つの時刻に絞り込む．
投票には組織内コミュニケーションツールを利用し，今回はSlackで実装した．
これらはユーザの自由意思の尊重，心理的負担の低減，および個人のプライバシーの保護を設計の軸としている．


実証実験により滞在履歴に大きな変化は見られなかったものの，アンケートにおける回答者の3割以上が本システムを実際の帰宅行動に活用していたと確認された．
これは一部帰宅時のコミュニケーションの発生に寄与していると言える．
以上より，帰宅タイミングに着目した同行支援というアプローチ自体は一定の有効性を有していると考えられる．
一方で，システムの意図や利用方法の伝達，およびユーザへの働きかけの強度には課題が残っており，これらの点については今後の改善が必要である．


\section{今後の課題}
今後の課題として主にアプローチ対象の拡張と提示するタイミングや提示時刻の調整，UIの改良，実証実験期間の追加の4点が挙げられる．

%アプローチ対象の話
% 本研究ではバスの時刻と組み合わせてアプローチを行ったため，他の公共交通機関を利用するユーザや，車などのより柔軟に帰宅時刻を変更できるユーザへのアプローチを検討する必要がある．
% 加えて，本研究では帰宅時のみに絞り支援を検討したが退室時全体を対象としたアプローチの検討も必要である．
本研究では，バスの時刻と組み合わせた帰宅同行支援を行ったため，公共交通機関の中でもバスを利用する被験者を主な対象としていた．
そのため，鉄道など他の公共交通機関を利用するユーザや，車など帰宅時刻を比較的柔軟に変更できるユーザに対しては，本システムの効果が十分に発揮されない可能性がある．
今後は，バスと鉄道など複数の公共交通機関を組み合わせて時刻を提示する方法や，公共交通機関に依存しない形で帰宅タイミングの近さを示す指標の導入など，移動手段の多様性に対応したアプローチの検討が必要である．
また，本研究では帰宅時のみに着目して支援を行ったが，退室時全体を対象としたアプローチについても検討の余地がある．



% 時刻提示開始のタイミングが時刻算出のタイミングに依存してるねって話いる
% これは滞在者のみにアプローチする視点からの設計だけど
% 近づいた判定する幅広げて提示開始のタイミングを一意に決めて統一してしまうパターンもできたね
% 提示時刻も毎回クラスタリングする影響で同じ人物・日付でもちょっとずつ誤差があったね
% 厳密に決める必要はないけどDBに蓄積して平す（やりすぎだから書くかはびみょい）
% その日の


提示UIに関しては，時刻情報の提示と同時に，本システムの目的や利用方法が一目で分かる短い説明文を表示し，初見のユーザにも意図が伝わりやすくする必要があると考えられる．
長期的に運用する場合，初見のユーザに十分な説明が行われない状況も想定されるため，UI上で役割や使い方を補足する工夫が重要となる．
また，投票状況の表示が背景色と馴染み，変化が分かりにくい場合があったため，視認性の向上を含めたUIの改良が求められる．
これらの改善により，ユーザがシステムの状態を直感的に把握でき，参加しやすくなると期待される．


一方，Slack上でのDM通知および投票UIに関しては，投票時のフィードバックが十分でなかった点や，投票締切時刻が明示されていなかった点が課題として挙げられる．
今後は，投票完了時に視覚的なフィードバックを与える演出や，締切時刻を明示する表示の追加などにより，システムの状態を直感的に理解できるUI設計が求められる．
さらに，投票操作自体の負担軽減も重要な課題である．
例えば，数字入力だけでなくスタンプによる投票など，より直感的で容易な入力方式を採用すると，参加率の向上が期待される．
加えて，DMの送信対象についても検討の余地がある．
過度な通知はユーザが通知を見なくなる要因となるため避ける必要があるが，運用の結果，1日あたりの時刻提示回数は1〜3回程度にとどまっていた．
この程度の頻度であれば過剰な通知にはなりにくいと考えられるため，DMの送信対象を全滞在者に拡張すると投票の参加率の向上が期待される．
また，SlackのDM上において投票を促す簡潔な説明文や行動を促進する文言を付加し，ユーザが意図を理解しやすくなる工夫もできると考えられる．


%実験期間の話
さらに，システム稼働期間が4週間と短期間であったため一時的な変動の影響を受けやすく，安定的な変化を評価するには不十分であった．
安定的な行動変容や長期的な利用傾向を評価するためには，長期運用による継続的なデータ収集と分析を行う必要がある．

% Local Variables: 
% mode: japanese-LaTeX
% TeX-master: "root"
% End: 
